\subsection{serial}\label{serial}
Das \verb|serial| Modul kapselt die serielle Ein\-/Ausgabe über UART
(optional USB\-CDC). Es erfüllt zwei Aufgaben: Zum einen leitet es die
Standardausgabe um -- die Newlib\-Hooks \verb|_write| und \verb|_read| binden
\verb|printf| an den UART bzw.\ USB, sodass Debug\-Ausgaben direkt
hinausgehen. Zum anderen ist es selbst eine \verb|kybd_t|\-Implementierung
(Abschnitt \ref{keyboard}, Gerät \verb|serial_dev|), die empfangene Zeichen
als Tasteneingabe in einen \verb|state_t| überführt.

Das Modul verbindet mehrere andere Bausteine: \verb|buffer| und
\verb|buffer_pool| (Abschnitte \ref{buffer}, \ref{buffer_pool}) für die
RX\-/TX\-Puffer, \verb|_time| (Abschnitt \ref{time}) zur Durchsatzmessung
(\verb|srxhdl|, \verb|stxhdl|) und \verb|cycle| (Abschnitt \ref{cycle}) für
Zeitstempel.

\subsubsection*{Betriebsarten}
Das Verhalten wird über das Bitfeld \verb|print_e| (Listing \ref{print_e})
gesteuert und mit \verb|serial_mode_set()| gesetzt. Es wählt u.\,a.\ die
Transport\-Schicht (\verb|USE_UART|, \verb|USE_DMA_RX|, \verb|USE_USB|) und
optionale Aufbereitung wie Zeitstempel (\verb|TIMESTAMP|), Lücken\-Erkennung
(\verb|GAP_DETECT|) oder reine Rohausgabe (\verb|RAW|).

\begin{listing}
\begin{minted}{C}
typedef enum {
    RAW        = 1,
    TIMESTAMP  = 2,
    GAP_DETECT = 4,
    ONE_SHOT   = 0x8,
    USE_UART   = 0x10,
    USE_DMA_RX = 0x20,
    USE_DMA_TX = 0x40, // funktioniert nicht
    USE_USB    = 0x80,
    MEASURE_BYTE_PER_SECONDS = 0x100,
} print_e;
\end{minted}
\caption{Ausgabe\-Betriebsarten \texttt{print\_e}}
\label{print_e}
\end{listing}

\subsubsection*{Zustand}
Die Konfiguration wird beim Initialisieren als \verb|sio_t| (Listing
\ref{sio_t}) übergeben: UART\-Handle, RX\-/TX\-Puffer, ein \verb|cycle_t| und
die Betriebsart. Den Laufzeitzustand hält das Modul intern in der einzigen
Variablen \verb|isio| (Typ \verb|isio_t|); diese ergänzt u.\,a.\ den
Gerätezustand \verb|state_t|, einen \verb|buffer_pool_t| sowie Zähler für
Überläufe und verworfene USB\-Bytes.

\begin{listing}
\begin{minted}{C}
typedef struct _sio_t {
    UART_HandleTypeDef *uart;
    buffer_t *buffer[SIO_RXTX_CNT]; // [SIO_RX], [SIO_TX]
    cycle_t  *cycle;
    print_e   mode;
} sio_t;
\end{minted}
\caption{Konfigurations\-Struktur \texttt{sio\_t}}
\label{sio_t}
\end{listing}

\subsubsection*{Empfang}
{\sloppy\emergencystretch=3em
Der Empfang läuft per DMA mit Idle\-Line\-Erkennung
(\verb|HAL_UARTEx_ReceiveToIdle_DMA|). Bei Empfang ruft die HAL den
\verb|HAL_UARTEx_RxEventCallback| auf: dieser kopiert die Daten in den
RX\-Puffer, wandelt sie in Kleinbuchstaben (\verb|buffer_tolower|) und gibt
sie über \verb|serial_apply_change()| an den \verb|state_t| weiter -- eine als
Hexzahl erkannte Eingabe schaltet die entsprechende Taste.
\par}

\subsubsection*{API}
{\sloppy\emergencystretch=3em
\begin{description}
\item[\texttt{serial\_mode\_set(mode)} / \texttt{serial\_mode\_get()}] Setzt
  bzw.\ liest das \verb|print_e|\-Bitfeld.
\item[\texttt{serial\_reset(dev)}] Setzt RX\-/TX\-Puffer und den Gerätezustand
  zurück.
\item[\texttt{serial\_waitForNumber(key)}] Liefert die zuletzt empfangene
  Eingabe als Zahl bzw.\ $-1$; \verb|key| zeigt danach auf die Zeichenkette.
\item[\texttt{serial\_get\_byte\_per\_second()} /
  \texttt{serial\_reset\_byte\_per\_second()}] Lesen bzw.\ Zurücksetzen des
  gemessenen Sende\-Durchsatzes.
\item[\texttt{serial\_uart\_printf(line, len)}] Gibt \verb|len| Bytes direkt
  über den UART aus.
\end{description}
\par}

\noindent Über die exportierte \verb|kybd_t| \verb|serial_dev| ist das Modul
zugleich als Eingabegerät an der \verb|keyboard|\-Schnittstelle nutzbar; die
Struktur \verb|device_t serial_io| bietet daneben einen generischen
Open/Read/Write\-Zugang.
